
\documentclass{article}
\usepackage{colortbl}
\usepackage{makecell}
\usepackage{multirow}
\usepackage{supertabular}

\begin{document}

\newcounter{utterance}

\centering \large Interaction Transcript for game `cladder', experiment `full\_v1.5\_default', episode 5038 with llama{-}3.1{-}8b{-}instruct.
\vspace{24pt}

{ \footnotesize  \setcounter{utterance}{1}
\setlength{\tabcolsep}{0pt}
\begin{supertabular}{c@{$\;$}|p{.15\linewidth}@{}p{.15\linewidth}p{.15\linewidth}p{.15\linewidth}p{.15\linewidth}p{.15\linewidth}}
    \# & \multicolumn{2}{c}{Player} && \multicolumn{2}{c}{Game Master} \\
    \hline

    \theutterance \stepcounter{utterance}  
    & & & \multicolumn{4}{p{0.6\linewidth}}{
        \cellcolor[rgb]{0.9,0.9,0.9}{
            \makecell[{{p{\linewidth}}}]{
                \texttt{\tiny{[P1$\langle$GM]}}
                \texttt{You are an expert in causal inference. The following question is not a typical commonsense query, but rather a meticulously designed question created by a professor specializing in causal inference, intended to assess the students' mastery of the course content.} \\
\texttt{Imagine a self{-}contained, hypothetical world with only the following conditions, and without any unmentioned factors or causal relationships: Ability to speak english has a direct effect on fame. Talent has a direct effect on fame. The overall probability of speaking english is 89\%. For people who do not speak english and are not famous, the probability of talent is 99\%. For people who do not speak english and are famous, the probability of talent is 97\%. For people who speak english and are not famous, the probability of talent is 99\%. For people who speak english and are famous, the probability of talent is 57\%. If we look at people who are famous, does the chance of talent decrease when speaking english?} \\
\texttt{Start your answer with 'Yes' or 'No', followed by additional reasoning or evidence to support your explanation.} \\
            }
        }
    }
     \\ \\

    \theutterance \stepcounter{utterance}  
    & \multicolumn{4}{p{0.6\linewidth}}{
        \cellcolor[rgb]{0.561,0.737,0.561}{
            \makecell[{{p{\linewidth}}}]{
                \texttt{\tiny{[P1$\rangle$GM]}}
                \texttt{No, the chance of talent does not decrease when speaking english among people who are famous.} \\
\\ 
\texttt{To answer this question, we need to apply the concept of causal inference, specifically the backdoor adjustment and the frontdoor adjustment. However, in this case, we can use the frontdoor adjustment to directly compare the probability of talent between the two groups.} \\
\\ 
\texttt{Let's denote the variables as follows: } \\
\\ 
\texttt{{-} E: English speaker (E=1) or non{-}English speaker (E=0)} \\
\texttt{{-} F: Famous (F=1) or not famous (F=0)} \\
\texttt{{-} T: Talented (T=1) or not talented (T=0)} \\
\\ 
\texttt{We are given the following conditional probabilities:} \\
\\ 
\texttt{{-} P(T|E=0, F=0) = 0.99} \\
\texttt{{-} P(T|E=0, F=1) = 0.97} \\
\texttt{{-} P(T|E=1, F=0) = 0.99} \\
\texttt{{-} P(T|E=1, F=1) = 0.57} \\
\\ 
\texttt{We want to compare the probability of talent between the two groups of famous people: those who speak English and those who do not.} \\
\\ 
\texttt{To do this, we can use the law of total probability to calculate the probability of talent among famous people who speak English and those who do not:} \\
\\ 
\texttt{P(T|F=1) = P(T|E=1, F=1) * P(E=1|F=1) + P(T|E=0, F=1) * P(E=0|F=1)} \\
\\ 
\texttt{We know that P(E=1) = 0.89, so P(E=0|F=1) = 1 {-} P(E=1|F=1). We can use the law of total probability to calculate P(E=1|F=1):} \\
\\ 
\texttt{P(E=1|F=1) = P(E=1, F=1) / P(F=1)} \\
\\ 
\texttt{However, we are not given the joint probability P(E=1, F=1) directly. We can use the law of total probability to calculate it:} \\
\\ 
\texttt{P(E=1, F=1) = P(E=1) * P(F=1|E=1) + P(E=0) * P(F=1|E=0)} \\
\\ 
\texttt{We know that P(E=0) = 1 {-} P(E=1) = 0.11, and we can use the law of total probability to calculate P(F=1|E=1) and P(F=1|E=0):} \\
\\ 
\texttt{P(F=1|E=1) = P(F=1, E=1) / P(E=1)} \\
\texttt{P(F=1|E=0) = P(F=1, E=0) / P(E=0)} \\
\\ 
\texttt{However, we are not given the joint probabilities P(F=1, E=1) and P(F=1} \\
            }
        }
    }
    & & \\ \\

    \theutterance \stepcounter{utterance}  
    & & & \multicolumn{2}{p{0.3\linewidth}}{
        \cellcolor[rgb]{0.9,0.9,0.9}{
            \makecell[{{p{\linewidth}}}]{
                \texttt{\tiny{[GM$|$GM]}}
                \texttt{No} \\
            }
        }
    }
    & & \\ \\

    \theutterance \stepcounter{utterance}  
    & & & \multicolumn{2}{p{0.3\linewidth}}{
        \cellcolor[rgb]{0.9,0.9,0.9}{
            \makecell[{{p{\linewidth}}}]{
                \texttt{\tiny{[GM$|$GM]}}
                \texttt{yes} \\
            }
        }
    }
    & & \\ \\

    \theutterance \stepcounter{utterance}  
    & & & \multicolumn{2}{p{0.3\linewidth}}{
        \cellcolor[rgb]{0.9,0.9,0.9}{
            \makecell[{{p{\linewidth}}}]{
                \texttt{\tiny{[GM$|$GM]}}
                \texttt{game\_result = LOSE} \\
            }
        }
    }
    & & \\ \\

\end{supertabular}
}

\end{document}
